\section{Discussion}

\paragraph{Not another monolithic binary.}
WorldEpisode should have a small semantic core and multiple bindings. LeRobot can be the
policy-training view; Rerun the multirate physical-data view; NCore the sensor/reconstruction
capture view; MCAP/ROS the raw log view; OpenUSD/SimReady the composed simulation-world view; glTF
or SPZ the Gaussian asset view; and a reference directory the conformance profile.

\paragraph{Lineage-safe evaluation.}
Episode-level random splits can leak the same reconstructed room, physical object, Gaussian asset,
source video, world revision, or generated asset lineage across training and evaluation.
WorldEpisode should support split constraints such as \texttt{world\_lineage},
\texttt{physical\_entity}, \texttt{source\_capture}, \texttt{reconstruction\_run}, and
\texttt{generated\_asset\_family}.

\paragraph{Standards alignment.}
The project should align with standards bodies rather than antagonize them. Static 3D robot-map
exchange, humanoid dataset lifecycle, OpenLABEL annotation, glTF Gaussian splats, USD composition,
and NCore capture conventions are all useful neighbors. WorldEpisode is an executable profile and
bridge across those layers.

\paragraph{Reviewer red team.}
Likely criticisms are predictable: ``this is just USD,'' ``GSDF already combines splats and
physics,'' ``Rerun already solves robot data,'' or ``a schema is not research.'' The response is to
ship adapters, conformance fixtures, fault-injection evaluation, replay metrics, and lineage-split
evidence rather than relying on schema novelty alone.
